\documentclass[10pt,a4paper]{amsart}
\usepackage[margin=19mm]{geometry}
\usepackage{amsmath,amssymb,mathtools}
\usepackage[scr=rsfs,cal=euler]{mathalfa}
\usepackage{xfrac}
\usepackage[unicode,hidelinks,bookmarks=false]{hyperref}
\newcommand{\ie}{i.e.\ }
\newcommand{\cf}{cf.\ }
\newcommand{\resp}{respectively\ }
\newcommand{\Subsection}{\subsection}
\providecommand{\nobreakdash}{\nobreak\hbox{-}}
\setlength{\emergencystretch}{2em}
\multlinegap0pt
\usepackage{amssymb}
\usepackage{multirow}
\usepackage{tikz}
\usepackage{caption}
\newtheorem{lemma}{Lemma}
\newtheorem{claim}{Claim}
\newtheorem{proofofclaim}{Proofofclaim}
\usepackage{cleveref}
\crefname{lemma}{Lemma}{Lemmas}
\crefname{claim}{Claim}{Claims}
\crefname{proofofclaim}{Proofofclaim}{Proofofclaims}
\newcommand{\C}{\mathcal{C}}
\newcommand{\R}{\mathcal{R}}
\newcommand{\set}[1]{\left\{#1\right\}}
\title{Pushing Cops and Robber on Graphs of Maximum Degree 4}
\author{Harmender Gahlawat}
\date{Source: arXiv:2505.21450v2 (CC BY 4.0)}
\begin{document}
\maketitle

\noindent\emph{Source and licence.} Pushing Cops and Robber on Graphs of Maximum Degree 4. Version arXiv:2505.21450v2 (CC BY 4.0), distributed by arXiv under the Creative Commons Attribution 4.0 International licence, http://creativecommons.org/licenses/by/4.0/, as recorded in the arXiv licence metadata for this exact version and shown on the abstract page https://arxiv.org/abs/2505.21450v2. Full licence notice: \texttt{licenses/arxiv-2505.21450-CC-BY-4.0.txt}.\par\smallskip

\noindent\textit{Editorial scope.} Counted unit: the article's claim C:neighborvisited (source0.tex lines 482--484). The article's complete original proof of it is delivered with it (source0.tex lines 485--493, one \texttt{proof} environment of 9 lines). No mathematics is written, repaired or re-derived: the statement and the proof are the article's own lines, in the article's own notation, and no retained line is substituted or re-flowed.  Retained uncounted as the source-local context the proof needs: lines 355--357.  Nothing was omitted from any retained span, so the package carries no omission record.  No retained line contains a citation, so the wrapper carries no bibliography and no bibliography entry.  No retained line contains a figure environment, an image-inclusion command or drawing code, so no image is delivered and the package declares no asset dependency.  Editorial formatting only, in addition: an AMS \texttt{amsart} preamble that restates the article's own declarations and macros used by the retained lines, namely the environments \texttt{lemma}, \texttt{claim}, \texttt{proofofclaim} on separate counters, together with the standard packages those lines need.  The retained environments print in this document as Lemma 1, Claim 1, Proofofclaim 1 in order of appearance, whereas the article prints the same items with the numbering of its own full document.  The complete native source of the article as distributed in the arXiv e-print for this exact version is bundled unchanged as \texttt{sources/178870/source0.tex}.
\begin{lemma}\label{L:trivial}
    If $\R$ occupies a vertex $v\in V(\overrightarrow{G})$ such that $|N^+(v)| \leq 1$ on a cop's move, then $\R$ will be trapped in the next cop move.
\end{lemma}

    \begin{claim}\label{C:neighborvisited}
        Given $xy\notin E(G)$, if $|N^+(x)| \neq 2$ or $|N^+(y)| \neq 2$, then $\C$ can trap $\R$ in at most $2$ robber moves.
    \end{claim}

    \begin{proofofclaim}
        Without loss of generality, let us assume that $|N^+(x)|\neq 2$. $\C$ begins with pushing $y$, forcing $\R$ to move to $x$. Notice that this does not change $|N^+(x)|$ since $xy\notin E(G)$. We distinguish the following two cases:

        \medskip
        \noindent\textbf{Case~1.} $|N^+(x)|\leq 1$: $\C$ can trap $\R$ using \cref{L:trivial}.

        \medskip
        \noindent\textbf{Case~2.} $|N^+(x)| =3$. $\C$ pushes $x$, and now $N^+(x) = \set{u}$. This forces $\R$ to move to $u$. Since we have pushed both out-neighbors of $u$, $N^+(u) = \emptyset$, and thus $\R$ is trapped. 
    \end{proofofclaim}

\end{document}
